%% kapa: exam=1 round=24 year=2013 period=1 form=A subject=economics subject_ko=경제원론 seq=20 number=60
%% source: Q-Net 감정평가사 기출문제 (cst00309) artlSeq=5073802 file=extracted/1008/5073802/2013년 제24회 감정평가사 1차시험 기출문제/1교시(A형).hwp
%% conversion: kordoc 4.24.51 -> tools/gen1.py (auto); review=proofread-sonnet; flags=none
\begin{question}{60}
외부효과가 존재하는 A시장의 수요곡선은 $P=100-Q$이고, 사적한계비용은 $PMC=40+0.5Q$이다. 생산량 한 단위당 30의 추가적인 사회적 비용이 발생하는 경우에 관한 설명으로 옳은 것은? (단, $P$ 는 가격, $Q$ 는 수량이다.)
\begin{choices}
  \item 정부개입이 없는 경우 균형생산량은 20이다.
  \item 사회적 후생을 극대화하는 생산량은 40이다.
  \item 보조금을 지급하여 사회적 후생을 높일 수 있다.
  \item 생산량 수준을 40으로 규제함으로써 사회적 후생을 높일 수 있다.
  \item 생산량 수준을 20으로 규제하든 단위당 30의 조세를 부과하든 사회적 후생의 크기는 동일하다.
\end{choices}
\end{question}
\answer{3,5}
\answernote{복수 정답(정답표 원문: 3, 5)}
